% Mục 4 (bản rút gọn). Rút từ sections/corr_*.tex. Nguồn: tests/*.cpp, mk/part{1,2,3}.mk,
% results/{apple-m5,C4M7Y670NT,LAP1184}/{test,tsan,mutation}.log.
\subsection{Bài kiểm tra}
Nhóm kiểm tính an toàn (không có hai thread cùng ở trong critical section, không thread nào ra khỏi
barrier sớm, không giá trị counter nào bị trả về hai lần) và tính sống còn (chương trình không treo).
Công bằng chỉ được đo. Cùng một bài kiểm tra dạng template chạy cho mọi cài đặt, kể cả cài đặt của thư
viện. Nhóm dùng bốn kỹ thuật:
\begin{itemize}[nosep]
  \item stress test với ít thread (1--4), với số luồng phần cứng và với gấp đôi số đó; ở mức gấp đôi,
    thread bị tước core giữa giao thức;
  \item watchdog, thread giám sát báo \texttt{FAIL} khi bài kiểm tra ngừng tiến triển;
  \item ThreadSanitizer (TSan), công cụ của compiler báo hai truy cập xung đột, ít nhất một không atomic,
    mà không có quan hệ happens-before;
  \item mutation testing, chạy bộ kiểm tra trên các biến thể cố ý làm hỏng.
\end{itemize}

Với lock, critical section tăng một biến đếm không atomic và đếm số thread đang ở bên trong, nên vi phạm
hiện ra trực tiếp hoặc qua một lần cập nhật bị mất. Lượt đầu chạy số vòng cố định, lượt sau chạy trong
300~ms. Watchdog báo deadlock sau 20~s. Với barrier, ở episode $e$ mỗi thread ghi $e$ vào ô của mình, gọi
\texttt{wait}, rồi đọc mọi ô. Một ô nhỏ hơn $e$ nghĩa là thread này qua sớm. Phép kiểm handoff dùng xen
kẽ hai barrier với các ô là biến thường, nên dưới TSan nó kiểm thêm memory order. Bản schedule fuzzing
chèn một độ trễ ngẫu nhiên vào mỗi điểm giữa hai bước của giao thức.
% nguồn: tests/test_locks.cpp:220-232 (ms = 300, watchdog 20 s); tests/test_barrier.cpp.

Với counter, nhóm lưu mọi giá trị trả về và kiểm năm điều: giá trị cuối bằng $M$ (\emph{final}), không
giá trị nào bị trả về hai lần (\emph{dup}) hay bị thiếu (\emph{miss}), dãy giá trị của mỗi thread tăng
ngặt (\emph{order}), và thứ tự thời gian thực (\emph{real-time}). Phép kiểm sau cùng kẹp mỗi lời gọi giữa
hai lần \texttt{fetch\_add} \texttt{seq\_cst} lên một đồng hồ logic dùng chung. Nếu A trả về trước khi B
bắt đầu thì A phải nhận giá trị nhỏ hơn. Các phép kiểm này kiểm
linearizability~\cite{herlihy1990linearizability} trên các lịch chạy đã xảy ra. Ở \texttt{-O2},
compiler gộp một vòng tăng không atomic thành một lần đọc và một lần ghi. Trong một thí nghiệm phụ, bài
kiểm tra chỉ so giá trị cuối không phát hiện lỗi của một counter không đồng bộ ở 198/200 lần chạy với 2
thread. Vì vậy bài kiểm tra của nhóm cho địa chỉ counter đi qua \texttt{pc::opaque()} ở mỗi vòng.
% nguồn: tests/test_counter.cpp:278-295, 340-411; docs/part3.md mục 5.3 (2 thread có start gate sai 2/200).

Trên cả ba máy, mọi cài đặt đạt mọi bài kiểm tra, và TSan không báo data race nào trên mã đúng. Bản TSan
chạy với kích thước nhỏ hơn 10--20 lần. Ở bản thường, MCS ở gấp đôi số luồng phần cứng chỉ lấy lock
15\,221 lần trong 300~ms trên \PCname{apple-m5}, so với 2\,705\,564 lần ở 4 thread, phù hợp với việc lock
được trao cho một thread đang không có core (RQ2). \PCname{lock_mcs} hết ngân sách thời gian và chỉ được
kiểm qua khoảng 2 triệu lần tăng trên mỗi máy (1\,765\,312 lần trên \PCname{apple-m5}). TSan không kiểm
được \PCname{omp_lock} và \PCname{omp_barrier}, vì libgomp không được instrument, nên nhóm bỏ qua chúng ở
bản TSan. Vì cùng lý do với libc++ của macOS, bản TSan trên hai máy Apple bỏ phép kiểm handoff của
\PCname{std_barrier}.
% nguồn: results/apple-m5/test.log (dòng mcs; lock_mcs "time budget reached"); C4M7Y670NT 2\,049\,423,
% LAP1184 2\,063\,205 lần tăng; tsan.log (SKIP omp_lock, omp_barrier skipped, note std_barrier).

\subsection{Mutation testing}\label{sec:corr:mutation}
\begin{rqanswer}\textbf{Trả lời RQ1.} Không bài kiểm tra nào tìm thấy lỗi trên ba máy, và bộ kiểm tra bắt
cả 14 biến thể hỏng. Tính đúng dựa trên lập luận ở mục Cài đặt. Kiểm thử củng cố lập luận đó. Trên
\PCname{apple-m5}, \PCname{C4M7Y670NT} và \PCname{LAP1184}, mọi cài đặt đạt mọi bài kiểm tra (12 lock,
7 barrier và 19 counter trên hai máy Apple, 13, 9 và 20 trên \PCname{LAP1184}), và TSan không báo data
race nào. Hai lỗi memory order khó bắt nhất. Stress test không bắt \texttt{TAS\_RELEASE\_RELAXED} lần nào
(0/3 trên cả ba máy). Bản thường bắt \texttt{fetch\_add\_release\_only} 1/5 lần trên \PCname{apple-m5},
4/5 trên \PCname{C4M7Y670NT} và 0/5 trên \PCname{LAP1184}. TSan bắt cả hai lỗi trên mọi máy.\end{rqanswer}
% nguồn: results/{apple-m5,C4M7Y670NT,LAP1184}/mutation.log.

\begin{table}[tbp]\centering\scriptsize\setlength{\tabcolsep}{3pt}
% nguồn: results/apple-m5/mutation.log
\begin{tabularx}{\linewidth}{@{}>{\ttfamily}l>{\raggedright\arraybackslash}Xlll@{}}
\toprule
\normalfont\textbf{Biến thể} & \textbf{Lỗi cài vào} & \textbf{Stress (+ fuzzing)} & \textbf{TSan} & \textbf{Bị bắt bởi} \\
\midrule
TAS\_CHECK\_THEN\_SET & \texttt{load} rồi \texttt{store} thay \texttt{exchange} & 3/3 vi phạm loại trừ & race & cả hai \\
TAS\_RELEASE\_RELAXED & nhả lock bằng \texttt{store} relaxed & 0/3 & race & chỉ TSan \\
MCS\_NO\_WAIT\_SUCCESSOR & nhả lock không chờ thread kế tiếp & 3/3 treo & không & stress \\
MCS\_ALWAYS\_CLEAR\_TAIL & luôn đặt \texttt{tail} về null & 3/3 treo & race & cả hai \\
\midrule
single\_flag & một cờ, không đảo chiều & 5/5 (5/5) sai episode & không & stress \\
flip\_before\_reset & đảo \texttt{sense} trước khi đặt lại bộ đếm & 2/5 (5/5) treo & không & stress \\
nonatomic\_count & \texttt{load}+\texttt{store} thay \texttt{fetch\_add} & 5/5 (5/5) treo & không & stress \\
fetch\_add\_release\_only & \texttt{fetch\_add} release thay acq\_rel & 1/5 (0/5) sai episode & race & chủ yếu TSan \\
thread\_local\_sense & local sense chung cho mọi barrier & 5/5 (5/5) sai giá trị handoff & race & cả hai \\
\midrule
plain\_long & \texttt{long} thường & 3/3 final, dup, miss, order & race & cả hai \\
load\_store & \texttt{load} rồi \texttt{store} atomic & 3/3 final, dup, miss, order & không & kiểm giá trị \\
cas\_no\_retry & CAS một lần, không thử lại & 3/3 final, dup, miss & không & kiểm giá trị \\
lock\_read\_outside & đọc giá trị trước \texttt{acquire()} & 3/3 final, dup, miss, order & race & cả hai \\
block\_reserve & giữ trước 8 giá trị mỗi thread & 3/3 chỉ real-time & không & kiểm giá trị \\
\bottomrule
\end{tabularx}
\caption{Mutation testing trên \PCname{apple-m5}: lock (trên), barrier (giữa), counter (dưới). Cột stress
cho số lần chạy bắt được trên số lần chạy (trong ngoặc: bản schedule fuzzing) cùng dấu hiệu hoặc phép kiểm
báo lỗi; cột TSan cho biết TSan có báo data race không.}
\label{tab:corr:mutation}
\end{table}

\emph{Bằng chứng.} Bảng~\ref{tab:corr:mutation} tóm tắt kết quả trên \PCname{apple-m5}. Trên hai máy còn
lại, mọi biến thể đều bị bắt, chỉ khác ở số lần stress test bắt được. Riêng trên \PCname{LAP1184}, TSan
không báo \texttt{thread\_local\_sense}, nên biến thể đó chỉ bị phép kiểm giá trị bắt, và bản thường bắt
\texttt{flip\_before\_reset} 5/5 lần.
% nguồn: results/C4M7Y670NT/mutation.log, results/LAP1184/mutation.log.

\emph{Vì sao.} Hai biến thể chỉ TSan bắt ổn định đều là lỗi memory order. \texttt{TAS\_RELEASE\_RELAXED}
đổi lệnh nhả lock \texttt{stlrb} thành \texttt{strb}, và store relaxed không tạo cạnh happens-before giữa
hai chủ lock liên tiếp. \texttt{fetch\_add\_release\_only} đổi \texttt{ldaddal} thành \texttt{ldaddl}, nên
thread tới cuối có thể đọc ô của thread khác trước khi \texttt{fetch\_add} của nó có hiệu lực. Trên
x86-64, store relaxed và release đều là \texttt{mov}, và mọi \texttt{fetch\_add} đều là \texttt{lock
xadd}, nên trên \PCname{LAP1184} hai biến thể này chạy đúng lệnh máy của bản đúng. Ngược lại, TSan im
lặng với bảy biến thể chỉ dùng atomic hoặc gây treo, và stress test hoặc phép kiểm giá trị bắt mọi lỗi
thuật toán. Fuzzing làm \texttt{flip\_before\_reset} treo sau 3--7 episode, trong khi bản thường cần
39\,896 và 289\,052 episode ở hai lần bắt được trên \PCname{apple-m5}.
% nguồn: results/apple-m5/mutation_logs/barrier_mut2.run{2,5}.log, barrier_mut2_fuzz.run{1..5}.log.

\emph{Độ vững.} Stress test và fuzzing chỉ xét các lịch chạy đã xảy ra, nên chỉ ra được là có lỗi, không
chứng minh được là không có lỗi. TSan chỉ tìm data race, chỉ thấy các đường đi đã chạy và chỉ thấy mã được
instrument. Bộ kiểm tra không kiểm tiến triển (lock-free, wait-free) ngoài trường hợp treo, và các biến
thể hỏng do chính nhóm chọn. Mỗi máy chạy bộ kiểm tra một lượt, nên tỉ lệ bắt lỗi (ví dụ 2/5 hay 1/5)
chỉ cho bậc độ lớn. Trên \PCname{LAP1184}, ở lượt TSan 24 thread không lock nào lấy được lock lần nào,
và ở lượt 12 thread \PCname{mcs} và \PCname{ttas_bo_spin} cũng không. TSan vì vậy chỉ thực sự kiểm lượt
4 thread với mọi lock, và lượt 12 thread với các lock còn lại.
% nguồn: results/LAP1184/tsan.log, cột 24t = 0 acq ở cả 12 lock; cột 12t: mcs 0 acq, ttas_bo_spin 0 acq.
